% ----------------------------------------------------------------------
%             Plantilla basada en la de la Universidad of Deusto:
%             http://softwarelibre.deusto.es/wp-content/uploads/2011/
%             11/phd-template_v.1.0.zip
% ----------------------------------------------------------------------

\documentclass[twoside,12pt]{Latex/Classes/PhDthesisPSnPDF}
%---------------------------------------------------------------
% Macros
% version 3 by Igor Ruiz-Agundez 2011
% version 2 by Jakob Suckale 2007
% version 1 by Harish Bhanderi 2002
%---------------------------------------------------------------

% This file contains macros that can be called up from connected TeX files
% It helps to summarise repeated code, e.g. figure insertion (see below).



%---------------------------------------------------------------
% Figures
%---------------------------------------------------------------


% Makes the \InsertFig macro compatible both with one or two columns
\makeatletter
\newlength \figwidth
\if@twocolumn
  \setlength \figwidth {\columnwidth}
\else
  \setlength \figwidth {\textwidth}
\fi
\makeatother

% \InsertFig allows inserting figures
% Parameters
% 1 --> Filename
% 2 --> Label for referencing
% 3 --> Title describing the figure (caption)
% 4 --> Description of the figure
% 5 --> Figure width, range [0,1]. If parameter is left blank the figure size is not change
% 6 --> Any other option for \includegraphics
% Usage:
% \InsertFig{}{}{}{}{}{}
%
\newcommand{\InsertFig}[6]{%
	\ifthenelse{\isempty{#5}}%
	{% if #1 is empty
		\begin{figure}[htbp!]
		\centering
		\includegraphics[#6]{#1}%
		\caption{#3}{\textbf{#4}}
		\label{#2}
		\end{figure}    
	}
	{% if #1 is not empty
		\begin{figure}[htbp!]
		\centering
		\includegraphics[width=#5\figwidth,#6]{#1}%
		\caption{#3}{\textbf{#4}}
		\label{#2}
		\end{figure}
	}
}

%% Simple version of \InsertFig
%\newcommand{\InsertFig}[5]{
%  \begin{figure}[htbp]
%   	\centering
%    \includegraphics[width=#4\textwidth,#5]{#1}%
%    \caption{#3}
%    \label{#2}
%  \end{figure}
%}



% insert a centered figure with caption
% parameters 1:filename, 2:label, 3:title, 
\newcommand{\figuremacro}[3]{
	\begin{figure}[htbp]
		\centering
		\includegraphics[width=1\textwidth]{#1}
		\caption[#3]{\textbf{#3}}
		\label{#2}
	\end{figure}
}


% insert a centered figure with caption and description
% parameters 1:filename, 2:label, 3:title, 4:description
\newcommand{\figuremacroD}[4]{
	\begin{figure}[htbp]
		\centering
		\includegraphics[width=1\textwidth]{#1}
		\caption[#3]{\textbf{#3} - #4}
		\label{#2}
	\end{figure}
}

% insert a centered figure with caption and description AND WIDTH
% parameters 1:filename, 2:label, 3:title, 4:description, 5: textwidth
% textwidth 1 means as text, 0.5 means half the width of the text
\newcommand{\figuremacroDW}[5]{
	\begin{figure}[htbp]
		\centering
		\includegraphics[width=#5\textwidth]{#1}
		\caption[#3]{\textbf{#3} - #4}
		\label{#2}
	\end{figure}
}

% inserts a figure with wrapped around text; only suitable for NARROW figs
% o is for outside on a double paged document; others: l, r, i(inside)
% text and figure will each be half of the document width
% note: long captions often crash with adjacent content; take care
% in general: above 2 macro produce more reliable layout
\newcommand{\figuremacroN}[3]{
	\begin{wrapfigure}{o}{0.5\textwidth}
		\centering
		\includegraphics[width=0.48\textwidth]{#1}
		\caption[#2]{{\small\textbf{#2} - #3}}
		\label{#1}
	\end{wrapfigure}
}




% Estas definiciones son para el comando \InsertFigBox
\newlength{\anchoFigura}
\newlength{\anchoFloat}
\addtolength{\fboxsep}{2\fboxsep}
%\renewcommand{\capfont}{\normalfont\normalcolor\sffamily\small}
%\renewcommand{\caplabelfont}{\normalfont\normalcolor\sffamily\bfseries\small}

% El comando \InsertFigBox nos permite insertar figuras en un marco
% Los parametros son:
% 1 --> Fichero de la imagen
% 2 --> Etiqueta (label) para referencias
% 3 --> Texto a pie de imagen
% 4 -> Porcentaje del ancho de página que ocupará la figura (de 0 a 1)
% 5 --> Opciones que queramos pasarle al \includegraphics
\newcommand{\InsertFigBox}[5]{%
  \setlength{\anchoFloat}{#4\textwidth}%
  \addtolength{\anchoFloat}{-4\fboxsep}%
  \setlength{\anchoFigura}{\anchoFloat}%
  \begin{figure}%
    \begin{center}%
      \Ovalbox{%
        \begin{minipage}{\anchoFloat}%
          \begin{center}%
            \includegraphics[width=\anchoFigura,#5]{#1}%
            \caption{#3}%
            \label{#2}%
          \end{center}%
        \end{minipage}
      }%
    \end{center}%
  \end{figure}%
}



%---------------------------------------------------------------
% Misc
%---------------------------------------------------------------

% predefined commands by Harish
\newcommand{\PdfPsText}[2]{
  \ifpdf
     #1
  \else
     #2
  \fi
}


%---------------------------------------------------------------
% Locales
%---------------------------------------------------------------


%%
%% Para quitar traducciones raras (Cuadros)
%% A de usarse cada vez que se seleccione el idioma
%%
\newcommand{\MejorarTraducciones}{%
       \renewcommand{\listtablename}{Índice de tablas}
       \renewcommand{\tablename}{Tabla}
       \renewcommand{\lstlistingname}{Lista}
}%



%---------------------------------------------------------------
% Source code
%---------------------------------------------------------------


%%
%% Para escribir extractos de codigo
%%
%% Las tabulaciones se substituyen por dos espacios
%\fvset{tabsize=2}
%% Creamos un nuevo environment de fancyvrb para los ejemplos enmarcados
%\DefineVerbatimEnvironment{VerbEj}{BVerbatim}{fontsize=\small,samepage=true,commandchars=\\\{\}}
%% Colo de fondo
%\definecolor{grisfondo}{gray}{0.9}
%% Environment para extractos de codigo
%\newenvironment{codigo}%
%{\VerbatimEnvironment\begin{Sbox}\begin{VerbEj}}%
%{\end{VerbEj}\end{Sbox}\setlength{\fboxsep}{8pt}\begin{center}\fcolorbox{black}{grisfondo}{\TheSbox}\end{center}}
%
%% Otro formato más bonito para código fuente
%\newcommand{\codigofuente}[3]{%
%  \lstinputlisting[language=#1,caption={#2}]{#3}%
%}





\graphicspath{ {./figures/} }


\ifpdf  
    \pdfinfo { /Title  (PhD)
               /Creator (TeX)
               /Producer (pdfTeX)
               /Author (Rubén Héctor García Ortega)
               /CreationDate (D:20130705120000)
               /ModDate (D:YYYYMMDDhhmm)
               /Subject (xyz)
               /Keywords (keyword1, keyword2, keyword3) }
    \pdfcatalog { /PageMode (/UseOutlines)
                  /OpenAction (fitbh)  }
\fi


\title{MADE: Massive Artificial Drama Engine for non player characters}


\crest{\includegraphics[width=10cm]{logo_ugr.jpg}}
\university{}
\degree{Trabajo Fin de Master presentado por }
\author{Rubén Héctor García Ortega} 
\collegeordept{}
\textadvisor{Dirigido por }
\advisor{Dr. Juan Julián Merelo Güervós}
\advisortwo{}
\textsignaturecandidate{El alumno}
\textsignatureadvisor{El director}
\cityofbirth{Granada}
%\degreedate{\monthname \ \the\year}
\degreedate{Julio de \the\year}


\hbadness=10000
\hfuzz=50pt



\begin{document}

\selectlanguage{spanish}

% sets line spacing
\renewcommand\baselinestretch{1.2}
\baselineskip=18pt plus1pt



\maketitle

% Thesis Titleback ---------------------------------------------------

\thispagestyle{empty}

\hfill

\vfill

\medskip


\noindent
\textit{
MADE: Massive Artificial Drama Engine for non player characters
}




Autor: Rubén Héctor García Ortega

Director: Juan Julián Merelo Güervós




\vfill

\vfill

\noindent
Las siguientes páginas web contienen información actualizada sobre el presente trabajo: \\
\url{https://github.com/raiben/made/}\\
\url{https://www.velonuboso.com/made/}


\noindent
Impresión realizada en Granada

% TODO final date
\noindent
Primera edición, 
% Moth and year
\monthname \ \the\year

\vspace{1cm}
\hrule
\bigskip

% \cleardoublepage command ends the current page and causes all figures and tables that have so far appeared in the input to be printed. In a two-sided printing style, it also makes the next page a right-hand (odd-numbered) page, producing a blank page if necessary. 
\cleardoublepage

%%: ----------------------- cover page back side ------------------------
%% Your research institution may require reviewer names, etc.
%% This cover back side is required by Dresden Med Fac; uncomment if needed.
%
%\newpage
%\vspace{10mm}
%1. Reviewer: Name
%
%\vspace{10mm}
%2. Reviewer: 
%
%\vspace{20mm}
%Day of the defense:
%
%\vspace{20mm}
%\hspace{70mm}Signature from head of PhD committee:
%
%
%\cleardoublepage

% ----------------------------------------------------------------------

\frontmatter

% Thesis Dedictation ---------------------------------------------------

\begin{dedication} %this creates the heading for the dedication page

\textit{... a Olga y Óliver}

\end{dedication}

% ----------------------------------------------------------------------


% Thesis Abstract -----------------------------------------------------


%\begin{abstractslong}    %uncommenting this line, gives a different abstract heading


\begin{resumen}        %this creates the heading for the abstract page

\selectlanguage{spanish}

%Motivación



%Razonamiento

%Solución propuesta

%Resultados


\end{resumen}




%\end{abstractlongs}


% ---------------------------------------------------------------------- 


%% Thesis Acknowledgements ------------------------------------------------


% Opening of the acknowledgements

%Sort version
%this creates the heading for the acknowlegments
\begin{acknowledgements}      
%Long version
%uncommenting this line, gives a different acknowledgements heading
%\begin{acknowledgementslong} 

Lorem ipsum dolor sit amet, consectetur adipiscing elit. Ut ultrices egestas nunc, venenatis rhoncus elit fermentum non. Pellentesque gravida nulla vitae ipsum lobortis ullamcorper. Ut adipiscing, tellus in egestas mattis, enim metus pretium erat, ac tempor dolor neque placerat nulla. Nullam nec ligula eu ipsum pharetra semper a in magna. Integer ut tortor quis nisi fringilla euismod eu ac ipsum. Pellentesque sodales consectetur erat eget rutrum. Proin ornare dolor ut arcu aliquet vestibulum. Pellentesque laoreet tincidunt sem eget semper.

Integer interdum mattis magna ullamcorper tristique. Nullam commodo nulla eget ipsum vulputate tincidunt auctor leo aliquet. Fusce euismod sagittis ante, eu vulputate eros dictum at. Cras non euismod nunc. Nullam velit diam, consectetur sed eleifend vitae, blandit at arcu. Maecenas ut urna nec turpis lobortis commodo. Aliquam aliquet turpis id massa viverra id sollicitudin est cursus. Sed a tortor non mauris cursus imperdiet.

Integer fermentum rutrum urna at vestibulum. Vivamus ullamcorper erat in sapien dignissim pellentesque. Integer convallis fringilla dictum. In bibendum lectus eu nulla pretium volutpat. Morbi hendrerit fringilla tortor, sed gravida neque lacinia a. In risus magna, hendrerit vitae cursus ac, vehicula at eros. Aenean quis ipsum sit amet leo vestibulum cursus.

Cras placerat mattis dui quis vehicula. Nulla sit amet metus nibh, at auctor enim. Quisque congue ultricies sapien in suscipit. Fusce vitae placerat ante. Praesent aliquet urna ac elit consequat nec mattis augue faucibus. Nunc et sapien vel felis mollis sodales. Aenean molestie nulla vestibulum nisi fringilla vel euismod dolor tristique. Aenean fermentum, dolor eget tincidunt faucibus, risus lorem feugiat elit, sagittis malesuada eros ligula in odio. Pellentesque ac libero lobortis justo bibendum laoreet. Cras egestas lorem eget ligula dignissim sollicitudin. Vestibulum sit amet augue ultrices erat faucibus vestibulum. Aenean tincidunt faucibus leo, nec auctor diam bibendum a. Sed varius, mauris in pellentesque scelerisque, nisl ligula viverra erat, in eleifend tellus enim ac magna. Pellentesque quis est risus. Cras mollis feugiat auctor. Proin ac eros vitae nulla gravida varius.

Morbi at augue sapien. Duis tempus quam vitae velit interdum ultricies. Vivamus laoreet lacinia elit sit amet vehicula. Ut congue diam ac magna hendrerit sed fermentum justo lacinia. Curabitur vel odio neque, quis consequat mi. Proin lobortis justo quis enim fermentum accumsan sagittis ipsum imperdiet. Proin sem felis, laoreet placerat egestas id, fringilla id mauris. Pellentesque a nisi sit amet leo consectetur gravida nec et dui. Curabitur quis hendrerit augue. Etiam sed dui nec tortor convallis fringilla. Proin tempor mattis diam nec egestas. Quisque condimentum elementum lacus ac porta. Vivamus congue, odio eu ullamcorper elementum, leo turpis tempus sem, at condimentum dolor quam eu nunc. Pellentesque eget risus ac velit aliquam sollicitudin sed et ipsum.

Donec nulla enim, scelerisque quis dignissim ut, vehicula sit amet felis. Ut tristique pulvinar aliquet. Proin vitae odio nibh. Sed est dolor, malesuada et lobortis vitae, tempor quis tellus. Proin mollis lacus eget arcu tempor vitae bibendum nisi adipiscing. In sollicitudin pretium dapibus. Sed accumsan imperdiet diam quis pellentesque.

Sed in lacinia lectus. Nullam ac magna in erat blandit posuere nec nec ante. Pellentesque habitant morbi tristique senectus et netus et malesuada fames ac turpis egestas. Maecenas ipsum neque, fringilla sed sodales id, scelerisque et erat. Morbi placerat mauris vel elit ultricies sed molestie ipsum faucibus. Nunc vel sem orci, ac placerat justo. Nulla varius nisi mauris. Vivamus sed ligula felis, ac mattis purus. Vivamus felis nibh, bibendum et aliquam sed, feugiat a elit. Morbi imperdiet libero vitae lacus viverra tempus. Praesent ultricies fermentum urna eget varius. Nunc blandit, augue eget vehicula venenatis, ligula eros posuere tellus, non sollicitudin tortor velit quis augue. Nulla rhoncus fringilla dolor, vel cursus ligula varius id. Vestibulum et eros eros, ut euismod elit. Praesent placerat libero sed lorem mollis sed cursus erat laoreet.

Suspendisse in fermentum nulla. Donec blandit ultricies felis, non suscipit turpis sagittis eget. Donec eget enim in erat auctor lobortis eu sagittis mauris. Aliquam vitae nibh nisi. Integer nulla erat, feugiat quis vulputate vel, vulputate id erat. Morbi dolor dui, volutpat congue fermentum tempus, rutrum et quam. Vivamus aliquam, sapien in ultrices dictum, tortor tellus dictum justo, eu viverra velit tellus congue nibh. Suspendisse potenti. Nullam at massa id leo sollicitudin ultricies ut et nibh. Vivamus fermentum pellentesque ante et lobortis. Etiam eros justo, rutrum at varius et, porttitor ac lorem. Aenean viverra, mauris eget imperdiet aliquet, risus quam tincidunt risus, quis elementum enim diam a ipsum. Ut et arcu nunc, a viverra mi. Nullam egestas lectus vel ipsum sodales eget consectetur metus facilisis. Proin egestas magna sed felis mattis ultrices.

Morbi a lectus vitae lacus malesuada tincidunt. Vestibulum eu quam justo. Integer hendrerit posuere augue, in fringilla nulla feugiat ut. Fusce vel libero sed tellus posuere molestie in sed justo. Nunc orci magna, aliquam id dictum sit amet, bibendum vestibulum augue. Maecenas luctus ultricies elit, et aliquet justo rutrum sed. Proin mauris dui, cursus quis tempus non, rhoncus nec ante. Lorem ipsum dolor sit amet, consectetur adipiscing elit.

Fusce ac nunc non felis fermentum ultrices. Curabitur nibh felis, convallis id consequat rhoncus, elementum sodales ipsum. In hac habitasse platea dictumst. In quis nulla eu sapien mattis cursus nec semper leo. Maecenas sagittis viverra quam, vel sodales elit placerat et. Sed condimentum ultricies mattis. Pellentesque scelerisque convallis lobortis. In tempus lorem nec lorem blandit sodales. 


\begin{flushright}
\textit{Eskerrik asko,}

Your name

% Moth and year
\monthname \ \the\year



% Signature figure

%\begin{figure}[htbp!]
%\end{figure}
%\includegraphics{signature}%



\end{flushright}



%Closing of the acknowledgements
%Sort version
\end{acknowledgements}
% Long version
%\end{acknowledgementslong}

% ------------------------------------------------------------------------




\selectlanguage{spanish}

\setcounter{secnumdepth}{5} % organisational level that receives a numbers
\setcounter{tocdepth}{5}    % print table of contents for level 3



\tableofcontents            % print the table of contents
\listoffigures	% print list of figures
\listoftables  % print list of tables


% this file is called up by thesis.tex
% content in this file will be fed into the main document

% Glossary entries are defined with the command \nomenclature{1}{2}
% 1 = Entry name, e.g. abbreviation; 2 = Explanation
% You can place all explanations in this separate file or declare them in the middle of the text. Either way they will be collected in the glossary.

% required to print nomenclature name to page header
\markboth{\MakeUppercase{\nomname}}{\MakeUppercase{\nomname}}

% ----------------------- contents from here ------------------------
%

%
%
%% acronyms


\nomenclature{IP}{Internet Protocol}







%\begin{multicols}{2} % \begin{multicols}{#columns}[header text][space]
%\begin{footnotesize} % scriptsize(7) < footnotesize(8) < small (9) < normal (10)
%\printnomenclature[1.5cm] % [] = distance between entry and description

\printnomenclature % [] = distance between entry and description

\label{sec:glossary} % target name for links to glossary

%\end{footnotesize}
%\end{multicols}




%: --------------------------------------------------------------
%:                  MAIN DOCUMENT SECTION
% --------------------------------------------------------------


\mainmatter

\renewcommand{\chaptername}{} % uncomment to print only "1" not "Chapter 1"
\pagestyle{fancy}

%: ----------------------- subdocuments ------------------------

% Parts of the thesis are included below. Rename the files as required.
% But take care that the paths match. You can also change the order of appearance by moving the include commands.

%: ----------------------- introduction ------------------------
% introduction

\begin{savequote}[50mm]
The beginning is the most important part of the work. 
\qauthor{Plato}
\end{savequote}

\chapter{Introduction}
\label{cha:Introduction}

% the code below specifies where the figures are stored
\ifpdf
    \graphicspath{{1_introduction/figures/PNG/}{1_introduction/figures/PDF/}{1_introduction/figures/}}
\else
    \graphicspath{{1_introduction/figures/EPS/}{1_introduction/figures/}}
\fi


%------------------------------------------------------------------------- 

\cite{turing1950computing}

Lorem ipsum dolor sit amet, consectetur adipiscing elit. Ut ultrices egestas nunc, venenatis rhoncus elit fermentum non. Pellentesque gravida nulla vitae ipsum lobortis ullamcorper. Ut adipiscing, tellus in egestas mattis, enim metus pretium erat, ac tempor dolor neque placerat nulla. Nullam nec ligula eu ipsum pharetra semper a in magna. Integer ut tortor quis nisi fringilla euismod eu ac ipsum. Pellentesque sodales consectetur erat eget rutrum. Proin ornare dolor ut arcu aliquet vestibulum. Pellentesque laoreet tincidunt sem eget semper.

Integer interdum mattis magna ullamcorper tristique. Nullam commodo nulla eget ipsum vulputate tincidunt auctor leo aliquet. Fusce euismod sagittis ante, eu vulputate eros dictum at. Cras non euismod nunc. Nullam velit diam, consectetur sed eleifend vitae, blandit at arcu. Maecenas ut urna nec turpis lobortis commodo. Aliquam aliquet turpis id massa viverra id sollicitudin est cursus. Sed a tortor non mauris cursus imperdiet.

Integer fermentum rutrum urna at vestibulum. Vivamus ullamcorper erat in sapien dignissim pellentesque. Integer convallis fringilla dictum. In bibendum lectus eu nulla pretium volutpat. Morbi hendrerit fringilla tortor, sed gravida neque lacinia a. In risus magna, hendrerit vitae cursus ac, vehicula at eros. Aenean quis ipsum sit amet leo vestibulum cursus.

Cras placerat mattis dui quis vehicula. Nulla sit amet metus nibh, at auctor enim. Quisque congue ultricies sapien in suscipit. Fusce vitae placerat ante. Praesent aliquet urna ac elit consequat nec mattis augue faucibus. Nunc et sapien vel felis mollis sodales. Aenean molestie nulla vestibulum nisi fringilla vel euismod dolor tristique. Aenean fermentum, dolor eget tincidunt faucibus, risus lorem feugiat elit, sagittis malesuada eros ligula in odio. Pellentesque ac libero lobortis justo bibendum laoreet. Cras egestas lorem eget ligula dignissim sollicitudin. Vestibulum sit amet augue ultrices erat faucibus vestibulum. Aenean tincidunt faucibus leo, nec auctor diam bibendum a. Sed varius, mauris in pellentesque scelerisque, nisl ligula viverra erat, in eleifend tellus enim ac magna. Pellentesque quis est risus. Cras mollis feugiat auctor. Proin ac eros vitae nulla gravida varius.

Morbi at augue sapien. Duis tempus quam vitae velit interdum ultricies. Vivamus laoreet lacinia elit sit amet vehicula. Ut congue diam ac magna hendrerit sed fermentum justo lacinia. Curabitur vel odio neque, quis consequat mi. Proin lobortis justo quis enim fermentum accumsan sagittis ipsum imperdiet. Proin sem felis, laoreet placerat egestas id, fringilla id mauris. Pellentesque a nisi sit amet leo consectetur gravida nec et dui. Curabitur quis hendrerit augue. Etiam sed dui nec tortor convallis fringilla. Proin tempor mattis diam nec egestas. Quisque condimentum elementum lacus ac porta. Vivamus congue, odio eu ullamcorper elementum, leo turpis tempus sem, at condimentum dolor quam eu nunc. Pellentesque eget risus ac velit aliquam sollicitudin sed et ipsum. 



% ----------------------------------------------------------------------

	

\section{Motivación}
\label{sec:Motivation}

%
%\InsertFig{computing_machinery_and_intelligence}{fig:Computing machinery and intelligence}{Computing machinery and intelligence.}{See Macro.tex for a detailed explanation of the InsertFig function}{1}{}
%


Lorem ipsum dolor sit amet, consectetur adipiscing elit. Ut ultrices egestas nunc, venenatis rhoncus elit fermentum non. Pellentesque gravida nulla vitae ipsum lobortis ullamcorper. Ut adipiscing, tellus in egestas mattis, enim metus pretium erat, ac tempor dolor neque placerat nulla. Nullam nec ligula eu ipsum pharetra semper a in magna. Integer ut tortor quis nisi fringilla euismod eu ac ipsum. Pellentesque sodales consectetur erat eget rutrum. Proin ornare dolor ut arcu aliquet vestibulum. Pellentesque laoreet tincidunt sem eget semper.

Integer interdum mattis magna ullamcorper tristique. Nullam commodo nulla eget ipsum vulputate tincidunt auctor leo aliquet. Fusce euismod sagittis ante, eu vulputate eros dictum at. Cras non euismod nunc. Nullam velit diam, consectetur sed eleifend vitae, blandit at arcu. Maecenas ut urna nec turpis lobortis commodo. Aliquam aliquet turpis id massa viverra id sollicitudin est cursus. Sed a tortor non mauris cursus imperdiet.

Integer fermentum rutrum urna at vestibulum. Vivamus ullamcorper erat in sapien dignissim pellentesque. Integer convallis fringilla dictum. In bibendum lectus eu nulla pretium volutpat. Morbi hendrerit fringilla tortor, sed gravida neque lacinia a. In risus magna, hendrerit vitae cursus ac, vehicula at eros. Aenean quis ipsum sit amet leo vestibulum cursus.


% ----------------------------------------------------------------------


	

%: ----------------------- related work ------------------------
% introduction

%\begin{savequote}[50mm]
%Personally, I think it does help, that it makes a beneficial difference, but the scientific literature on the subject is very messy.
%\qauthor{Jeanne Petrek}
%%“And upon the top of the pillars was lily work: so was the work of the pillars finished.”
%%
%% Bible quotes
%\end{savequote}


\chapter{Estado del arte}
\label{cha:State of the Art}


Lorem ipsum dolor sit amet, consectetur adipiscing elit. Ut ultrices egestas nunc, venenatis rhoncus elit fermentum non. Pellentesque gravida nulla vitae ipsum lobortis ullamcorper. Ut adipiscing, tellus in egestas mattis, enim metus pretium erat, ac tempor dolor neque placerat nulla. Nullam nec ligula eu ipsum pharetra semper a in magna. Integer ut tortor quis nisi fringilla euismod eu ac ipsum. Pellentesque sodales consectetur erat eget rutrum. Proin ornare dolor ut arcu aliquet vestibulum. Pellentesque laoreet tincidunt sem eget semper.

Integer interdum mattis magna ullamcorper tristique. Nullam commodo nulla eget ipsum vulputate tincidunt auctor leo aliquet. Fusce euismod sagittis ante, eu vulputate eros dictum at. Cras non euismod nunc. Nullam velit diam, consectetur sed eleifend vitae, blandit at arcu. Maecenas ut urna nec turpis lobortis commodo. Aliquam aliquet turpis id massa viverra id sollicitudin est cursus. Sed a tortor non mauris cursus imperdiet.

Integer fermentum rutrum urna at vestibulum. Vivamus ullamcorper erat in sapien dignissim pellentesque. Integer convallis fringilla dictum. In bibendum lectus eu nulla pretium volutpat. Morbi hendrerit fringilla tortor, sed gravida neque lacinia a. In risus magna, hendrerit vitae cursus ac, vehicula at eros. Aenean quis ipsum sit amet leo vestibulum cursus.

Cras placerat mattis dui quis vehicula. Nulla sit amet metus nibh, at auctor enim. Quisque congue ultricies sapien in suscipit. Fusce vitae placerat ante. Praesent aliquet urna ac elit consequat nec mattis augue faucibus. Nunc et sapien vel felis mollis sodales. Aenean molestie nulla vestibulum nisi fringilla vel euismod dolor tristique. Aenean fermentum, dolor eget tincidunt faucibus, risus lorem feugiat elit, sagittis malesuada eros ligula in odio. Pellentesque ac libero lobortis justo bibendum laoreet. Cras egestas lorem eget ligula dignissim sollicitudin. Vestibulum sit amet augue ultrices erat faucibus vestibulum. Aenean tincidunt faucibus leo, nec auctor diam bibendum a. Sed varius, mauris in pellentesque scelerisque, nisl ligula viverra erat, in eleifend tellus enim ac magna. Pellentesque quis est risus. Cras mollis feugiat auctor. Proin ac eros vitae nulla gravida varius.

Morbi at augue sapien. Duis tempus quam vitae velit interdum ultricies. Vivamus laoreet lacinia elit sit amet vehicula. Ut congue diam ac magna hendrerit sed fermentum justo lacinia. Curabitur vel odio neque, quis consequat mi. Proin lobortis justo quis enim fermentum accumsan sagittis ipsum imperdiet. Proin sem felis, laoreet placerat egestas id, fringilla id mauris. Pellentesque a nisi sit amet leo consectetur gravida nec et dui. Curabitur quis hendrerit augue. Etiam sed dui nec tortor convallis fringilla. Proin tempor mattis diam nec egestas. Quisque condimentum elementum lacus ac porta. Vivamus congue, odio eu ullamcorper elementum, leo turpis tempus sem, at condimentum dolor quam eu nunc. Pellentesque eget risus ac velit aliquam sollicitudin sed et ipsum. 






% ----------------------------------------------------------------------

	

%: ----------------------- overall methodology ------------------------
% introduction


%\begin{savequote}[50mm]
%Historical methodology, as I see it, is a product of common sense applied to circumstances. 
%\qauthor{Samuel E. Morison}
%\end{savequote}


\chapter{Metodología general}
\label{cha:Overall methodology}

La metodología de evaluación del sistema MADE consitirá en el siguiente proceso:

\begin{itemize}
 \item Definición de los parámetros de ejecución del Entorno MADE
 \item Definición de la función fitness
 \item Estudio de la variabilidad del fitness para soluciones aleatorias
 \item Definición de los parámetros de ejecución del algoritmo genético
 \item Ejecución del algoritmo para diferente número de perfiles
 \item Obtención de la mejor solución para los diferentes perfiles
 \item Análisis de los valores obtenidos
\end{itemize}

\section{Definición de los parámetros de ejecución del Entorno MADE}

El entorno MADE se ejecutará con los siguientes parámetros:

\begin{verbatim}
global.NUMBER_OF_PROFILES=1
global.NUMBER_OF_INITIAL_AGENTS=30
global.MAP_DIMENSION=30
global.FOOD=20
global.DAYS=2000

base.BASE_DAYS = 200
base.BASE_ENERGY = 5
base.BASE_SMELL = 3
base.BASE_NUTRITION = 4
base.BASE_BITE = 5
base.BASE_FUR = 5
base.BASE_AGE_TO_BE_ADULT_FEMALE = 42
base.BASE_AGE_TO_BE_ADULT_MALE = 49
base.BASE_PREGNANCY_TIME = 30
\end{verbatim}

\section{Definición de la función Fitness}





\section{Estudio de la variabilidad del fitness}

Se estudiará la variabilidad del fitness para 100 cromosomas, mediante la
ejecución de 100 Entornos MADE utilizando dichos cromosomas.

Utilizando los datos anteriores, se elegirá un número de iteraciones que
ofrezca una media significativa, y se utilizará para calcular el fitness medio
de cada indivíduo de la población del algoritmo genético.


\section{Parámetros de ejecución del algoritmo genético}

\begin{verbatim}
 global.MAX_ALLOWED_EVOLUTIONS=30
 global.POPULATION_SIZE=20
 global.NUMBER_OF_PROFILES={1,2,3,4}
\end{verbatim}

	
%: ----------------------- experiments and results ------------------------
% introduction

% this file is called up by thesis.tex
% content in this file will be fed into the main document

%: ----------------------- introduction file header -----------------------
\begin{savequote}[50mm]
The logic of validation allows us to move between the two limits of dogmatism and skepticism. 
\qauthor{Paul Ricoeur}
\end{savequote}


\chapter{Validation of the methodology}
\label{cha:Validation of the methodology}


%------------------------------------------------------------------------- 

\cite{turing1950computing}

Lorem ipsum dolor sit amet, consectetur adipiscing elit. Ut ultrices egestas nunc, venenatis rhoncus elit fermentum non. Pellentesque gravida nulla vitae ipsum lobortis ullamcorper. Ut adipiscing, tellus in egestas mattis, enim metus pretium erat, ac tempor dolor neque placerat nulla. Nullam nec ligula eu ipsum pharetra semper a in magna. Integer ut tortor quis nisi fringilla euismod eu ac ipsum. Pellentesque sodales consectetur erat eget rutrum. Proin ornare dolor ut arcu aliquet vestibulum. Pellentesque laoreet tincidunt sem eget semper.

Integer interdum mattis magna ullamcorper tristique. Nullam commodo nulla eget ipsum vulputate tincidunt auctor leo aliquet. Fusce euismod sagittis ante, eu vulputate eros dictum at. Cras non euismod nunc. Nullam velit diam, consectetur sed eleifend vitae, blandit at arcu. Maecenas ut urna nec turpis lobortis commodo. Aliquam aliquet turpis id massa viverra id sollicitudin est cursus. Sed a tortor non mauris cursus imperdiet.

Integer fermentum rutrum urna at vestibulum. Vivamus ullamcorper erat in sapien dignissim pellentesque. Integer convallis fringilla dictum. In bibendum lectus eu nulla pretium volutpat. Morbi hendrerit fringilla tortor, sed gravida neque lacinia a. In risus magna, hendrerit vitae cursus ac, vehicula at eros. Aenean quis ipsum sit amet leo vestibulum cursus.

Cras placerat mattis dui quis vehicula. Nulla sit amet metus nibh, at auctor enim. Quisque congue ultricies sapien in suscipit. Fusce vitae placerat ante. Praesent aliquet urna ac elit consequat nec mattis augue faucibus. Nunc et sapien vel felis mollis sodales. Aenean molestie nulla vestibulum nisi fringilla vel euismod dolor tristique. Aenean fermentum, dolor eget tincidunt faucibus, risus lorem feugiat elit, sagittis malesuada eros ligula in odio. Pellentesque ac libero lobortis justo bibendum laoreet. Cras egestas lorem eget ligula dignissim sollicitudin. Vestibulum sit amet augue ultrices erat faucibus vestibulum. Aenean tincidunt faucibus leo, nec auctor diam bibendum a. Sed varius, mauris in pellentesque scelerisque, nisl ligula viverra erat, in eleifend tellus enim ac magna. Pellentesque quis est risus. Cras mollis feugiat auctor. Proin ac eros vitae nulla gravida varius.

Morbi at augue sapien. Duis tempus quam vitae velit interdum ultricies. Vivamus laoreet lacinia elit sit amet vehicula. Ut congue diam ac magna hendrerit sed fermentum justo lacinia. Curabitur vel odio neque, quis consequat mi. Proin lobortis justo quis enim fermentum accumsan sagittis ipsum imperdiet. Proin sem felis, laoreet placerat egestas id, fringilla id mauris. Pellentesque a nisi sit amet leo consectetur gravida nec et dui. Curabitur quis hendrerit augue. Etiam sed dui nec tortor convallis fringilla. Proin tempor mattis diam nec egestas. Quisque condimentum elementum lacus ac porta. Vivamus congue, odio eu ullamcorper elementum, leo turpis tempus sem, at condimentum dolor quam eu nunc. Pellentesque eget risus ac velit aliquam sollicitudin sed et ipsum. 










% ----------------------------------------------------------------------

	
	
%: ----------------------- conclusion ------------------------
% introduction

% this file is called up by thesis.tex
% content in this file will be fed into the main document

%: ----------------------- introduction file header -----------------------
%\begin{savequote}[50mm]
%Now this is not the end. It is not even the beginning of the end. But it is, perhaps, the end of the beginning. 
%\qauthor{Winston Churchill}
%\end{savequote}


\chapter{Conclusiones}
\label{cha:Conclusions}

Lorem ipsum dolor sit amet, consectetur adipiscing elit. Ut ultrices egestas nunc, venenatis rhoncus elit fermentum non. Pellentesque gravida nulla vitae ipsum lobortis ullamcorper. Ut adipiscing, tellus in egestas mattis, enim metus pretium erat, ac tempor dolor neque placerat nulla. Nullam nec ligula eu ipsum pharetra semper a in magna. Integer ut tortor quis nisi fringilla euismod eu ac ipsum. Pellentesque sodales consectetur erat eget rutrum. Proin ornare dolor ut arcu aliquet vestibulum. Pellentesque laoreet tincidunt sem eget semper.

Integer interdum mattis magna ullamcorper tristique. Nullam commodo nulla eget ipsum vulputate tincidunt auctor leo aliquet. Fusce euismod sagittis ante, eu vulputate eros dictum at. Cras non euismod nunc. Nullam velit diam, consectetur sed eleifend vitae, blandit at arcu. Maecenas ut urna nec turpis lobortis commodo. Aliquam aliquet turpis id massa viverra id sollicitudin est cursus. Sed a tortor non mauris cursus imperdiet.

Integer fermentum rutrum urna at vestibulum. Vivamus ullamcorper erat in sapien dignissim pellentesque. Integer convallis fringilla dictum. In bibendum lectus eu nulla pretium volutpat. Morbi hendrerit fringilla tortor, sed gravida neque lacinia a. In risus magna, hendrerit vitae cursus ac, vehicula at eros. Aenean quis ipsum sit amet leo vestibulum cursus.

Cras placerat mattis dui quis vehicula. Nulla sit amet metus nibh, at auctor enim. Quisque congue ultricies sapien in suscipit. Fusce vitae placerat ante. Praesent aliquet urna ac elit consequat nec mattis augue faucibus. Nunc et sapien vel felis mollis sodales. Aenean molestie nulla vestibulum nisi fringilla vel euismod dolor tristique. Aenean fermentum, dolor eget tincidunt faucibus, risus lorem feugiat elit, sagittis malesuada eros ligula in odio. Pellentesque ac libero lobortis justo bibendum laoreet. Cras egestas lorem eget ligula dignissim sollicitudin. Vestibulum sit amet augue ultrices erat faucibus vestibulum. Aenean tincidunt faucibus leo, nec auctor diam bibendum a. Sed varius, mauris in pellentesque scelerisque, nisl ligula viverra erat, in eleifend tellus enim ac magna. Pellentesque quis est risus. Cras mollis feugiat auctor. Proin ac eros vitae nulla gravida varius.

Morbi at augue sapien. Duis tempus quam vitae velit interdum ultricies. Vivamus laoreet lacinia elit sit amet vehicula. Ut congue diam ac magna hendrerit sed fermentum justo lacinia. Curabitur vel odio neque, quis consequat mi. Proin lobortis justo quis enim fermentum accumsan sagittis ipsum imperdiet. Proin sem felis, laoreet placerat egestas id, fringilla id mauris. Pellentesque a nisi sit amet leo consectetur gravida nec et dui. Curabitur quis hendrerit augue. Etiam sed dui nec tortor convallis fringilla. Proin tempor mattis diam nec egestas. Quisque condimentum elementum lacus ac porta. Vivamus congue, odio eu ullamcorper elementum, leo turpis tempus sem, at condimentum dolor quam eu nunc. Pellentesque eget risus ac velit aliquam sollicitudin sed et ipsum. 

% ----------------------------------------------------------------------


	
	

% --------------------------------------------------------------
%:                  BACK MATTER: appendices, refs,..
% --------------------------------------------------------------

% the back matter: appendix and references close the thesis
\backmatter


%: ----------------------- appendix ------------------------

%\appendix


%: ----------------------- bibliography ------------------------



% The section below defines how references are listed and formatted
% The default below is 2 columns, small font, complete author names.
% Entries are also linked back to the page number in the text and to external URL if provided in the BibTex file.


% Original version:

% PhDbiblio-url2 = names small caps, title bold & hyperlinked, link to page 
%\begin{multicols}{2} % \begin{multicols}{ # columns}[ header text][ space]
%\begin{tiny} % tiny(5) < scriptsize(7) < footnotesize(8) < small (9)
%
%\bibliographystyle{Latex/Classes/PhDbiblio-url2} % Title is link if provided
%\renewcommand{\bibname}{References} % changes the header; default: Bibliography
%
%\bibliography{9_backmatter/references} % adjust this to fit your BibTex file
%
%\end{tiny}
%\end{multicols}



% Show all bibliography entries
%\nocite*



% If we want bibliography backreference, use unsrt first and the desidered one after

%\bibliographystyle{unsrt} % Defines the bibliography style

\bibliographystyle{alpha} % Defines the bibliography style

%\bibliographystyle{apa-good} % Defines the bibliography style
%\bibliographystyle{natbib} % Defines the bibliography style

%\bibliographystyle{plainurl}

%\renewcommand{\bibname}{References} % changes the header; default: Bibliography

%To include the references/works cited/bibliography in your Table of Contents, right before the bibliography command, use the command
%\addcontentsline{toc}{section}{References}
%

\bibliography{references} % adjust this to fit your BibTex file

% --------------------------------------------------------------
% Various bibliography styles exit. Replace above style as desired.

% in-text refs: (1) (1; 2)
% ref list: alphabetical; author(s) in small caps; initials last name; page(s)
%\bibliographystyle{Latex/Classes/PhDbiblio-case} % title forced lower case
%\bibliographystyle{Latex/Classes/PhDbiblio-bold} % title as in bibtex but bold
%\bibliographystyle{Latex/Classes/PhDbiblio-url} % bold + www link if provided

%\bibliographystyle{Latex/Classes/jmb} % calls style file jmb.bst
% in-text refs: author (year) without brackets
% ref list: alphabetical; author(s) in normal font; last name, initials; page(s)

%\bibliographystyle{plainnat} % calls style file plainnat.bst
% in-text refs: author (year) without brackets
% (this works with package natbib)


% --------------------------------------------------------------


%: Declaration of originality
%
% this file is called up by thesis.tex
% content in this file will be fed into the main document

%: ----------------------- introduction file header -----------------------


% the code below specifies where the figures are stored
\ifpdf
    \graphicspath{{9_backmatter/figures/PNG/}{9_backmatter/figures/PDF/}{9_backmatter/figures/}}
\else
    \graphicspath{{9_backmatter/figures/EPS/}{9_backmatter/figures/}}
\fi

% Thesis statement of originality -------------------------------------

% Depending on the regulations of your faculty you may need a declaration like the one below. This specific one is from the medical faculty of the university of Dresden.

\begin{declaration}        %this creates the heading for the declaration page

I herewith declare that I have produced this work without the prohibited assistance of third parties and without making use of aids other than those specified; notions taken over directly or indirectly from other sources have been identified as such. This work has not previously been presented in identical or similar form to any examination board.

The dissertation work was conducted from 20XX to \the\year \ under the supervision of
Name Surname
and 
Name Surname
at 
the University of Deusto.

\vspace{10mm}

% City
Bilbao,

% Signature figure

%\begin{flushright}
%\begin{figure}[htbp!]
%\end{figure}
%\includegraphics{signature}%
%\end{flushright}


\end{declaration}



%-------------------------------------------------------------------------

\thispagestyle{empty}

\hfill
\vfill
\medskip

\begin{center}
\noindent
This dissertation was finished writing 
in 
Bilbao
on
% Date
\today 
\end{center}

\vfill
\medskip
\vspace{1cm}
\bigskip


%-------------------------------------------------------------------------

% Blank page
\clearpage
\thispagestyle{empty}

\hfill
\vfill
\medskip

\begin{center}
\textit{This page is intentionally left blank}
\end{center}

\vfill
\medskip
\vspace{1cm}
\bigskip


% ----------------------------------------------------------------------


\end{document}
